% Limitations and Future Work. Built by the future pass from the instructions
% below. Do not process now.

[Write Limitations and Future Work as honest prose with one supporting table. The
thesis depends on honest VVUQ, so state every gap plainly and do not approximate or
omit any of the steps that could not be run.]

[Subsection 1, limitations and approximations. Source this in full from the
\texttt{Limitations, approximations, and what could not be run} section of
\texttt{papers/VVUQ-02/execution/README.md}, and do not shorten it. Cover each
item: the on-prem LLM backends (anthropic, ollama, openai) were not installed and
no key was present, so the autonomy intent module
(\texttt{papers/VVUQ-02/codegen/src/autonomy/llm\_intent.py}) and the comparison
judge (\texttt{papers/VVUQ-02/codegen/src/llm/compare\_agent.py}) ran on their
deterministic offline paths; live Zenodo deposition was not performed, so
\texttt{papers/VVUQ-02/codegen/src/zenodo/patch\_pointers.py} wrote pointer JSON
with placeholder hashes rather than depositing the L0 stream; there was no mujoco,
NVIDIA Isaac, ROS 2, GPU, or display, so the behavior models are analytic
references and no physics-backed or rendered simulation ran; the committed sensor
sample is only the first 50 ms of the 60-second timeline (the full stream is the
Zenodo L0, by design not committed); local execution used a single Python (3.11.15)
with cross-version coverage delegated to the CI matrix; the paper templates were
intentionally not processed; and the H2-Surgical 1.0 is a hypothetical 2030 device
so every number is a simulation result and the four-entrant comparison is
simulation against simulation. Present these as a limitation-and-mitigation table,
abbreviating the module leaf names and giving the
\texttt{papers/VVUQ-02/codegen/src/} parent path once at the top.]

[Subsection 2, future work. Argue that the tractable next step is to adapt this
ten-gate, standards-anchored humanoid VVUQ automation across all oncology clinical
trial deliverables, because the speed, quality, and cost benefits shown here
generalise: each new deliverable inherits a reproducible determinism floor, a
pass-fail verification gate, validation against independent references, a
first-class uncertainty bound, and a runtime binding to published standards.
Sketch the path from this simulation-only reference toward physics-backed and
live-backend runs (mujoco or NVIDIA Isaac, on-prem LLM intent and judging, real
Zenodo deposition), then toward the regulatory artifacts a real deployment would
require (IEC 80601-2-77, IEC 60601, ISO 13482, FDA SaMD Class III clearance, IRB
approval, and authorization), citing \texttt{iec8060127}, \texttt{iec60601},
\texttt{iso13482}, and \texttt{fda2023credibility}. Connect to the broader
autonomous oncology-trial program (\texttt{kawchak2026sponsor},
\texttt{kawchak2026federated}) and frame the opportunity in terms of US leadership
in surgical-humanoid VVUQ for patient safety, efficacy, and trial speed.]
